If $\mu$ is a signed measure with $\mu 1$ finite and if
$$\nu = \lim_n [\mu(I + P + \cdots + P^{n-1})]$$
exists and is finite-valued, then $\mu$ is called a \text{left charge} with \text{potential measure} $\nu$. If $f$ is a function with $\alpha f$ finite and if
$$g = \lim_n [(I + P + \cdots + P^{n-1})f]$$
exists and is finite-valued, then $f$ is called a \text{right charge} with \text{potential function} $g$. In either case a \text{pure potential} is a potential of a non-negative charge.
